% ==============================================================================
% Section 1: Valuation Methodology & Theoretical Foundations
% ==============================================================================

\section[Valuation Methodology]{Valuation Methodology: Net Asset Value (Intrinsic Value) Method}

\subsection{Theoretical Framework and Valuation Premise}
In corporate financial accounting and quantitative equity valuation, the \textbf{Net Asset Value (NAV) Method}---also formally referred to as the \textbf{Liquidation Value Method} or \textbf{Asset-Backing Method}---determines the intrinsic equity value of an enterprise strictly through the prism of its tangible balance sheet resources. 

Under the classical accounting doctrine, the intrinsic value of a firm represents the residual financial claim that equity shareholders would legally receive if all company assets were liquidated at fair market value and all prior creditor obligations, debt securities, and external liabilities were settled in full.

\begin{definition}{Net Asset Value (Intrinsic Value)}{nav_def}
The \textbf{Net Asset Value} available to equity shareholders is defined as the aggregate fair value of all tangible assets less total external obligations. The resulting figure, when normalized across the total count of outstanding common shares, constitutes the \textbf{Intrinsic Value per Share} ($\IV$):
\begin{equation}
    \IV = \frac{\text{Net Assets Available to Equity Shareholders}}{\text{Number of Equity Shares Outstanding}}
\end{equation}
\end{definition}

\subsection{Mathematical Formulation}

\subsubsection{Goodwill Estimation via the Purchase of Average Profits Method}
Goodwill represents the capitalized value of super-normal earning power. Under the canonical textbook method, goodwill is valued at \textbf{three years' purchase of the average annual net profits} realized over the preceding five accounting periods:
\begin{equation}
    \overline{\Pi} = \frac{1}{5} \sum_{t=1}^{5} \Pi_t
\end{equation}
where $\Pi_t$ denotes the profit after tax for fiscal year $t$. The capitalized goodwill is subsequently evaluated as:
\begin{equation}
    \text{Goodwill} = \overline{\Pi} \times N_{\text{purchase}} = \overline{\Pi} \times 3
\end{equation}

\subsubsection{Net Assets Available to Equity Shareholders}
In empirical equity research, observable secondary market liquidation values for plant, property, equipment, and inventories are rarely disclosed in quarterly balance sheets. Consequently, audited balance sheet book values are adopted as the standard regulatory proxy. 

To maintain strict accounting conservatism and isolate the liquidation floor, existing unamortized or recorded goodwill is \textbf{excluded} from total assets, ensuring that only tangible resources back the equity:
\begin{equation}
    \text{Tangible Assets} = \text{Total Assets} - \text{Goodwill}
\end{equation}
\begin{equation}
    \text{Net Assets} = \text{Tangible Assets} - \text{Total Liabilities}
\end{equation}
Equivalently written in a single unified formulation:
\begin{equation}
    \text{Net Assets} = \text{Total Assets at Market/Book Value} - \text{Goodwill} - \text{Total Liabilities}
\end{equation}

\subsubsection{Intrinsic Value per Share}
Dividing the aggregate net tangible assets by the total volume of paid-up equity shares yields the per-share intrinsic worth:
\begin{equation}
    \boxed{\IV = \frac{\text{Net Assets}}{\text{Equity Shares Outstanding}}}
\end{equation}

\subsubsection{Market Divergence and Overvaluation Metrics}
To measure how the prevailing open-market quotation diverges from its underlying asset-backed intrinsic valuation, two formal metrics are defined:
\begin{enumerate}[label=\textbf{(\alph*)}]
    \item \textbf{Market Price per Share ($P_0$):}
    \begin{equation}
        P_0 = \frac{\text{Market Capitalization}}{\text{Number of Equity Shares Outstanding}}
    \end{equation}
    \item \textbf{Absolute Valuation Gap ($\Delta$):}
    \begin{equation}
        \Delta = \IV - P_0
    \end{equation}
    A negative valuation gap ($\Delta < 0$) indicates that the current market price trades at a substantial premium above the tangible asset base, signaling overvaluation under the liquidation framework.
    \item \textbf{Overvaluation Multiple / Ratio ($\mathcal{R}$):}
    \begin{equation}
        \mathcal{R} = \frac{P_0}{\IV} = \frac{\text{Market Capitalization}}{\text{Net Assets}}
    \end{equation}
    The ratio $\mathcal{R}$ quantifies the multiple of tangible net assets that investors in public capital markets are willing to pay for every unit of physical equity backing.
\end{enumerate}

\subsection{TikZ Process Flowchart: Net Asset Value Pipeline}
Figure~\ref{fig:nav_flowchart} details the operational filtration sequence of the Net Asset Value methodology from consolidated corporate assets down to intrinsic value and the comparative valuation gap.

\begin{figure}[htbp]
\centering
\begin{tikzpicture}[
    >=Stealth,
    node distance = 0.85cm and 1.0cm,
    every node/.style = {font=\small\sffamily},
    block/.style = {
        draw=pureblack,
        line width=1.1pt,
        fill=boxfill,
        rounded corners=2pt,
        text width=8.8cm,
        minimum height=1.1cm,
        align=center,
        inner sep=6pt
    },
    accentblock/.style = {
        draw=pureblack,
        line width=1.3pt,
        fill=tableheadbg,
        rounded corners=2pt,
        text width=8.8cm,
        minimum height=1.15cm,
        align=center,
        inner sep=6pt
    },
    resultblock/.style = {
        draw=pureblack,
        line width=1.4pt,
        fill=white,
        rounded corners=2pt,
        text width=8.8cm,
        minimum height=1.25cm,
        align=center,
        inner sep=6pt
    },
    conn/.style = {
        ->,
        line width=1.3pt,
        draw=pureblack
    }
]

    % Vertical Pipeline Nodes with strict symmetry and geometry
    \node [block] (n1) {
        \textbf{Consolidated Balance Sheet Assets}\\
        \footnotesize Aggregate Total Assets (Audited Book Value as Liquidation Proxy)
    };

    \node [block, below=0.85cm of n1] (n2) {
        \textbf{Less: Intangible Assets \& Goodwill}\\
        \footnotesize $\text{Tangible Assets} = \text{Total Assets} - \text{Recorded/Estimated Goodwill}$
    };

    \node [block, below=0.85cm of n2] (n3) {
        \textbf{Less: Total Liabilities \& Creditor Claims}\\
        \footnotesize Secured/Unsecured Borrowings, Provisions, and Current Liabilities
    };

    \node [accentblock, below=0.85cm of n3] (n4) {
        \textbf{Net Assets Available to Equity Shareholders}\\
        \footnotesize $\text{Net Assets} = \text{Tangible Assets} - \text{Total Liabilities}$
    };

    \node [accentblock, below=0.85cm of n4] (n5) {
        \textbf{Per-Share Normalization}\\
        \footnotesize $\text{Intrinsic Value per Share} = \frac{\text{Net Assets}}{\text{Number of Equity Shares Outstanding}}$
    };

    \node [resultblock, below=0.85cm of n5] (n6) {
        \textbf{Market Benchmarking \& Divergence Analysis}\\
        \footnotesize $\text{Valuation Gap} = \IV - P_0 \quad \Big| \quad \text{Overvaluation Multiple} = \frac{P_0}{\IV}$
    };

    % High-contrast connectors with clear shaft length (>= 0.85cm gap)
    \draw [conn] (n1) -- (n2);
    \draw [conn] (n2) -- (n3);
    \draw [conn] (n3) -- (n4);
    \draw [conn] (n4) -- (n5);
    \draw [conn] (n5) -- (n6);

\end{tikzpicture}
\caption{Net Asset Value (Intrinsic Value) Quantitative Filtration Pipeline.}
\label{fig:nav_flowchart}
\end{figure}
